\subsection{Preesquemas irreducibles y preesquemas conexos sobre un cuerpo
algebraicamente cerrado.}
\label{subsection:IV.2.4.4-fr}

\begin{env}[4.4.1]
\label{IV.2.4.4.1-fr}
Sean $k$ un cuerpo, $X$ un $k$-preesquema y $K$ una extensión de $k$. Como el
morfismo $\operatorname{Spec}(K)\to\operatorname{Spec}(k)$ es fielmente plano,
cuasicompacto y universalmente abierto (2.4.9), lo mismo ocurre con el morfismo
de proyección $p:X\otimes_k K\to X$ (2.2.13). Se deduce entonces de (2.3.5) que
toda componente irreducible de $X\otimes_k K$ \emph{domina} una componente
irreducible de $X$, y (como $p$ es sobreyectivo) toda componente irreducible de
$X$ está dominada por una componente irreducible de $X\otimes_k K$~; se obtiene
así de $p$ una aplicación \emph{sobreyectiva} del conjunto de las componentes
irreducibles de $X\otimes_k K$ en el conjunto de las componentes irreducibles
de $X$. Asimismo, del hecho de que $p$ es continuo y sobreyectivo se deduce que
toda componente conexa de $X$ es imagen por $p$ de una unión de componentes
conexas de $X\otimes_k K$, de donde resulta una aplicación \emph{sobreyectiva}
del conjunto de las componentes conexas de $X\otimes_k K$ en el conjunto de las
componentes conexas de $X$. Se concluye que toda componente irreducible (resp.
conexa) $Z'$ de $X\otimes_k K$ está contenida en un \emph{único} conjunto de la
forma $p^{-1}(Z)$, donde $Z$ es una componente irreducible (resp. conexa) de
$X$~; para todo subpreesquema de $X$ que tenga a $Z$ por espacio subyacente (y
que se denotará también por $Z$), $Z'$ es una componente irreducible (resp.
conexa) de $Z\otimes_k K$. En particular, si $X\otimes_k K$ tiene un número
finito $n$ (resp. $n'$) de componentes irreducibles (resp. conexas) (lo que
ocurrirá cuando $X$ sea \emph{de tipo finito} sobre $k$, pues entonces
$X\otimes_k K$ es de tipo finito sobre $K$ y, por tanto, noetheriano), el número
de componentes irreducibles (resp. conexas) de $X$ es $\leqslant n$ (resp.
$\leqslant n'$)~; decir que es igual a $n$ (resp. $n'$) significa que, para todo
subpreesquema reducido $Z$ de $X$ que tenga por espacio subyacente una
componente irreducible (resp. conexa) de $X$, $Z\otimes_k K$ es irreducible
(resp. conexo) (\textbf I, 5.1.8)~; el número de componentes irreducibles (resp.
conexas) de $X\otimes_k K$ es, por tanto, independiente de $K$ si y solo si,
para toda componente irreducible (resp. conexa) $Z$ de $X$, $Z\otimes_k K$ es
irreducible (resp. conexo) para \emph{toda} extensión $K$ de $k$.

En particular, si $X\otimes_k K$ es irreducible (resp. conexo), lo mismo ocurre
con $X$ (lo cual se deduce ya del hecho de que $p$ es sobreyectivo). Asimismo,
si $X\otimes_k K$ es \emph{reducido} (resp. \emph{íntegro}), lo mismo ocurre con
$X$, puesto que $p$ es fielmente plano (2.1.13).

En este número y en el siguiente examinaremos más de cerca qué puede decirse de
las componentes irreducibles (resp. conexas) de $X\otimes_k K$ cuando varía
$K$~; lo anterior nos muestra ya que el número de estas componentes es una función
\emph{creciente} de $K$.
\end{env}

\begin{lemma}[4.4.2]
\EGAReaderTarget{ega4.c04.lem.4-4-2}{IV.2.4.4.2-fr}
Sean $X'$, $X$ dos espacios topológicos y $f:X'\to X$ una aplicación continua.
Se supone que se satisfacen las condiciones siguientes~:
\begin{enumerate}
\item[\textup{(i)}] $f$ es sobreyectiva y abierta \textup{(resp.} tal que el espacio topológico
$X$ se identifica canónicamente con el cociente de $X$ por la relación de
equivalencia definida por $f$\textup{)}.

\oldpage[IV-2]{60}

\item[\textup{(ii)}] Para todo $x\in X$, $f^{-1}(x)$ es irreducible \textup{(resp.} conexo\textup{)}.
\end{enumerate}
Entonces $X'$ es irreducible \textup{(resp.} conexo\textup{)} si y solo si $X$ lo es.
\end{lemma}

\begin{proof}
La afirmación relativa a la conexidad no es sino Bourbaki,
\emph{Top. gén.}, chap.~I\textsuperscript{er}, 3\textsuperscript{e} éd.,
\S~11, n\textsuperscript{o}~3, prop.~7. Demostremos la afirmación relativa a
la irreducibilidad~; como la condición es trivialmente necesaria, puesto que
$f$ es sobreyectiva, demostremos que es suficiente. Sean, pues, $X'_1$, $X'_2$
dos partes cerradas de $X'$ tales que $X'=X'_1\cup X'_2$, y designemos por
$X_i$ $(i=1,2)$ el conjunto de los $x\in X$ tales que
$f^{-1}(x)\subset X'_i$~; se tiene entonces $X_i=X\setmin f(X'\setmin X'_i)$ y, como $f$ es
abierta, se concluye que $X_i$ es una parte \emph{cerrada} de $X$ $(i=1,2)$.
Por otra parte, para todo $x\in X$, $f^{-1}(x)$ es irreducible por hipótesis y
es unión de las dos partes cerradas $X'_1\cap f^{-1}(x)$ y
$X'_2\cap f^{-1}(x)$~; una de estas dos partes debe, por tanto, ser igual a
$f^{-1}(x)$, lo que significa que se debe tener $x\in X_1$ o $x\in X_2$~; se
tiene, pues, $X=X_1\cup X_2$ y, como $X$ es irreducible por hipótesis, se deduce
$X=X_1$ o $X=X_2$ y, por tanto, $X'=X'_1$ o $X'=X'_2$.
\end{proof}

\begin{remark}[4.4.3]
\label{IV.2.4.4.3-fr}
Si la topología de $X$ no es la topología de cociente de la de $X'$ por la
relación de equivalencia definida por $f$, puede ocurrir que $X$ y todas las
fibras $f^{-1}(x)$ sean irreducibles sin que $X'$ sea conexo~: se obtiene un
ejemplo tomando como $X$ un esquema irreducible (por ejemplo, la recta afín
$\operatorname{Spec}(k[T])$ sobre un cuerpo $k$ algebraicamente cerrado),
considerando un punto cerrado $x_0$ de $X$ y tomando como $X'$ la \emph{suma
topológica} de $\{x_0\}$ y del subespacio abierto $X\setmin\{x_0\}$ de $X$. Del mismo
modo, si se reemplaza la hipótesis «~$f$ abierta~» por «~$f$ cerrada~» (e
incluso «~$f$ propia~») (lo cual implica, sin embargo, que la topología de $X$
es entonces la topología de cociente de la de $X'$ por la relación definida
por $f$), puede ocurrir que $X$ y todas las fibras $f^{-1}(x)$ sean
irreducibles sin que $X'$ lo sea. Tomemos nuevamente como $X$ la recta afín
sobre $k$, designemos por $S$ el producto $X\times_k P$, donde
$P$ es la recta proyectiva sobre $k$~; si $x_0$ es un punto cerrado de
$X$, $t_0$ un punto cerrado de $P$, $p:S\to X$, $q:S\to P$ las
proyecciones, designemos por $X'$ el subpreesquema reducido cuyo espacio
subyacente es el conjunto cerrado $p^{-1}(x_0)\cup q^{-1}(t_0)$, y tomemos como
$f$ la restricción a $X'$ de $p$~; $f$ es propio, pero $X'$ no es irreducible.
\end{remark}

\begin{theorem}[4.4.4]
\label{IV.2.4.4.4-fr}
Sean $k$ un cuerpo algebraicamente cerrado y $X$ un $k$-preesquema. Si $X$ es
irreducible \textup{(resp.} conexo\textup{)}, lo mismo ocurre con $X\otimes_k K$ para toda
extensión $K$ de $k$.
\end{theorem}

\begin{proof}
En efecto, se ha visto (4.4.1) que el morfismo $p:X\otimes_k K\to X$ es
fielmente plano y abierto~; para poder aplicar el lema (4.4.2), basta verificar
que las fibras $p^{-1}(x)$ son irreducibles para todo $x\in X$. Pero el
$\boldsymbol{k}(x)$-preesquema $p^{-1}(x)$ es isomorfo a
$\operatorname{Spec}(\boldsymbol{k}(x)\otimes_k K)$ (\textbf I, 3.6.2) y es, por tanto, íntegro en
virtud de la hipótesis sobre $k$ y de (4.3.9).
\end{proof}

\begin{corollary}[4.4.5]
\label{IV.2.4.4.5-fr}
Sean $k$ un cuerpo algebraicamente cerrado, $X$ un $k$-preesquema, $K$ una
extensión de $k$ y $p:X\otimes_k K\to X$ la proyección canónica. Si $Z$ es una
parte irreducible \textup{(resp.} conexa\textup{)} de $X$, $p^{-1}(Z)$ es irreducible \textup{(resp.}
conexo\textup{)}~; en particular, si $X_0$ es una componente irreducible \textup{(resp.} la
componente conexa\textup{)} de $X$ que contiene a $Z$, $p^{-1}(X_0)$ es una componente
irreducible \textup{(resp.} conexa\textup{)} de $X\otimes_k K$ que contiene a $p^{-1}(Z)$.
\end{corollary}

\begin{proof}
La segunda afirmación se deduce de (4.4.1) y de la primera aplicada reemplazando
$Z$ por $X_0$~; para demostrar la primera, sean $X'$ el espacio subyacente a
$X\otimes_k K$ y $Z'=p^{-1}(Z)$~; la

\oldpage[IV-2]{61}

relación de equivalencia $R$ definida por $p$ sobre $X'$ es abierta, y lo mismo
ocurre, por tanto, con la relación $R_{Z'}$ que induce sobre la parte saturada
$Z'$; además, $Z$ se identifica con el espacio cociente $Z'/R_{Z'}$ (Bourbaki,
\emph{Top. gén.}, chap.~I\textsuperscript{er}, 3\textsuperscript{e} éd.,
\S~5, n\textsuperscript{o}~2, prop.~4)~; se puede aplicar, pues, a $Z$ y $Z'$ el
lema (4.4.2), de donde resulta la primera afirmación de (4.4.4).
\end{proof}

\begin{corollary}[4.4.6]
\label{IV.2.4.4.6-fr}
Con las hipótesis y notaciones de \textup{(4.4.5)}, la aplicación
$Z\leadsto p^{-1}(Z)$ es una biyección del conjunto de las componentes
irreducibles \textup{(resp.} conexas\textup{)} de $X$ sobre el conjunto de las componentes
irreducibles \textup{(resp.} conexas\textup{)} de $X\otimes_k K$, cuya biyección recíproca es
$Z'\leadsto p(Z')$.
\end{corollary}
